\chapter[Implementation]{Realization, implementation or experimentation}
\label{ch:implementation}

This chapter describes how the sprint-aligned design of
Section~\ref{sec:sprint-design} (Chapter~\ref{ch:architecture}) was realized in
code. Implementation follows the five-sprint plan of Table~\ref{tab:sprints},
with each subsection summarising the principal modules, technical choices, and
sequence flows validated for that increment.

\section{Hardware and Software Environment}

\begin{table}[H]
\centering
\caption{Development and deployment stack}
\label{tab:stack}
\begin{tabularx}{\textwidth}{>{\RaggedRight\arraybackslash}p{2.8cm} Y}
\toprule
Layer & Technology (versions) \\
\midrule
Workstation & Windows 10/11; Node.js 20+ \\
API & Express 4, TypeScript, Prisma 5, Vitest 3 \\
Web & Next.js 15, React 19, CSS Modules \\
Database & PostgreSQL 16 with pgvector extension \\
Queue & Redis 7, BullMQ 5 \\
AI services & Google Gemini (\tcode{gemini-2.5-flash}, \tcode{gemini-embedding-001}) \\
Containers & Docker Compose (dev and prod compose files) \\
\bottomrule
\end{tabularx}
\end{table}

The monorepo root orchestrates \texttt{npm run dev} (API on port 3001, web on 3000), database migrations, seeding, linting, and RAG evaluation scripts.

\section{Implementation Steps by Sprint}

\subsection{Sprint 1: Authentication and access control}
\label{subsec:impl-sprint1}

Sprint~1 implements the security perimeter defined in
Subsection~\ref{subsec:sprint1-design} (FR-01, FR-02, FR-12; NFR-02, NFR-03).
Access tokens are short-lived JWTs stored \textbf{in memory} on the web client---never
in \texttt{localStorage}---to reduce XSS exfiltration risk. Refresh tokens live in
HttpOnly cookies (\texttt{kp\_rt}). The client uses a \textbf{single in-flight}
refresh promise to prevent parallel rotations from invalidating each other~\cite{ietf-jwt-bcp}.
Role middleware and per-user restriction flags (\tcode{requireDocLibraryAccess},
\tcode{requireUseAiQueries}) enforce the hybrid authorization model of
Subsection~\ref{subsec:access-control-model}. Profile avatars are served from a
CORP-aware public route (\tcode{avatarsPublic.ts}). Flows
\ref{fig:seq-01}--\ref{fig:seq-07}, \ref{fig:seq-23}, \ref{fig:seq-32}, and
\ref{fig:seq-33} were validated against integration tests in
\tcode{auth.integration.test} and \tcode{avatars.integration.test}.

\subsection{Sprint 2: Document library and ingest pipeline}
\label{subsec:impl-sprint2}

Sprint~2 delivers the document lifecycle of Subsection~\ref{subsec:sprint2-design}
(FR-03 to FR-05). Uploads are validated (MIME allowlist, size limits), stored on
disk, and enqueued in BullMQ. The worker extracts text per format (PDF, Office,
HTML, etc.), chunks content with sentence-aware splitting and table preservation,
computes 768-dimensional embeddings, and persists \texttt{Chunk} rows with vector
indexes. Visibility rules (\texttt{ALL}, \texttt{DEPARTMENT}, \texttt{PRIVATE}) are
enforced through \tcode{canReadDocument} in \tcode{documentAccess.ts}, returning
404 on unauthorised access to prevent enumeration. Tags, favorites, recents, and
versioning are handled in \tcode{documents.ts}. Flows
\ref{fig:seq-08}--\ref{fig:seq-10} and \ref{fig:seq-24}--\ref{fig:seq-29} are
covered by \tcode{documents.integration.test}.

\subsection{Sprint 3: Search, RAG and conversations}
\label{subsec:impl-sprint3}

Sprint~3 implements the hybrid retrieval and AI assistant of
Subsection~\ref{subsec:sprint3-design} (FR-06 to FR-08; NFR-06).
Algorithm~\ref{alg:rag} summarizes the Ask endpoint logic in
\tcode{search.ts} and \tcode{ragCompletion.ts}.

\begin{algorithm}[H]
\caption{Hybrid RAG Ask pipeline (simplified)}
\label{alg:rag}
\begin{algorithmic}[1]
\State Authenticate user and resolve \texttt{readableDepartmentIds}
\State Optimize query (type, topic, rewrite, optional multi-hop)
\State Retrieve dense vectors (pgvector) and sparse FTS (BM25-style)
\State Fuse rankings with weighted reciprocal rank fusion (RRF)
\State Rerank top chunks with Gemini; optionally expand neighbor chunks
\State Stream SSE events: \texttt{sources}, \texttt{token}, optional \texttt{correction}
\State Persist conversation message with citations and confidence
\State Optionally cache answer keyed by prompt version and filters
\end{algorithmic}
\end{algorithm}

The Ask UI consumes SSE streams and renders markdown with sanitized links.
Conversation threads, message persistence, and thumbs-up/down feedback are
implemented in \tcode{conversations.ts}. Semantic search and Ask flows
(\ref{fig:seq-11}, \ref{fig:seq-12}) plus conversation sequences
(\ref{fig:seq-13}, \ref{fig:seq-14}, \ref{fig:seq-30}, \ref{fig:seq-31}) are
validated in \tcode{search.integration.test} and
\tcode{conversations.integration.test}.

\subsection{Sprint 4: Governance and notifications}
\label{subsec:impl-sprint4}

Sprint~4 covers administration, manager dashboards, and notifications per
Subsection~\ref{subsec:sprint4-design} (FR-09 to FR-11).
\texttt{departmentAccess.ts} computes readable and manageable department sets
including inherited manager/viewer scopes. Administrators manage users, bulk
restrictions, CSV import (with formula-injection escaping), department merge, KPI
exports, and document audit views in \tcode{admin.ts}. Managers access
department-scoped dashboards through \tcode{manager.ts}. The notification service
emits automatic events on document lifecycle changes and supports manual
announcements with attachments. Document previews use blob URLs for PDF/images;
notifications use polling with unread badges. Admin, manager, and notification
flows (\ref{fig:seq-15}--\ref{fig:seq-21}, \ref{fig:seq-34}--\ref{fig:seq-41})
are covered by \tcode{admin.integration.test}, \tcode{manager.integration.test},
and \tcode{notifications.integration.test}.

\subsection{Sprint 5: Operations, testing and deployment}
\label{subsec:impl-sprint5}

Sprint~5 consolidates operational readiness per
Subsection~\ref{subsec:sprint5-design} (NFR-01, NFR-04 to NFR-06). The health
probe (\ref{fig:seq-22}) reports PostgreSQL and Redis status via
\tcode{GET /health}. A 74-test Vitest/Supertest suite runs against real database
and Redis instances (\texttt{npm run test:api}). Docker Compose files provide
reproducible dev and production stacks with automatic migrations. The RAG
evaluation CLI (\texttt{npm run eval:rag:quick}, \texttt{npm run eval:rag})
writes JSON reports under \texttt{eval-reports/}. Table~\ref{tab:seq-index}
summarises all sprint sequence flows; the complete catalog appears in
Appendix~\ref{annex:diagrams} (Table~\ref{tab:annex-seq-index}).

\begin{table}[H]
\centering
\caption{Index of implementation sequence diagrams (Chapter~\ref{ch:architecture}, \S\ref{sec:sprint-design})}
\label{tab:seq-index}
\scriptsize
\begin{tabularx}{\textwidth}{|>{\Centering\arraybackslash}p{0.8cm}|>{\RaggedRight\arraybackslash}p{2.3cm}|>{\Centering\arraybackslash}p{1.0cm}|Y|}
\hline
\rowcolor{sectionblue!15}
\textbf{Sprint} & \textbf{Area} & \textbf{Fig.} & \textbf{Flow} \\
\hline
S1 & Auth & \ref{fig:seq-01}--\ref{fig:seq-07}, \ref{fig:seq-23}, \ref{fig:seq-32}, \ref{fig:seq-33} & Login, refresh, logout, password reset/change, profile, restrictions, /me, register info, restricted UX \\
\hline
S2 & Documents & \ref{fig:seq-08}--\ref{fig:seq-10}, \ref{fig:seq-24}--\ref{fig:seq-29} & Upload, ingest worker, download, list, detail, favorites, metadata, versions, reprocess \\
\hline
S3 & Search / RAG & \ref{fig:seq-11}, \ref{fig:seq-12} & Semantic search, hybrid Ask SSE \\
\hline
S3 & Conversations & \ref{fig:seq-13}, \ref{fig:seq-14}, \ref{fig:seq-30}, \ref{fig:seq-31} & Threads, messages, feedback, CRUD supplement, admin stats \\
\hline
S4 & Notifications & \ref{fig:seq-15}--\ref{fig:seq-17} & Automatic events, manual send, inbox \\
\hline
S4 & Manager & \ref{fig:seq-18} & Department dashboard \\
\hline
S4 & Admin & \ref{fig:seq-19}--\ref{fig:seq-21}, \ref{fig:seq-34}--\ref{fig:seq-41} & Users, departments, access, stats, audit, exports, bulk delete, lock \\
\hline
S5 & Operations & \ref{fig:seq-22} & Health check (tests, deploy, RAG eval: see \S\ref{subsec:impl-sprint5}) \\
\hline
\end{tabularx}
\end{table}

\section{Difficulties Encountered and Solutions}

\begin{table}[H]
\centering
\caption{Implementation challenges and resolutions}
\label{tab:difficulties}
\begin{tabularx}{\textwidth}{>{\RaggedRight\arraybackslash}p{2.8cm} >{\RaggedRight\arraybackslash}p{0.30\linewidth} Y}
\toprule
Difficulty & Analysis & Solution adopted \\
\midrule
Avatar not displayed & Browser blocked cross-origin image (CORP) & Set \tcode{Cross-Origin-Resource-Policy: cross-origin} on \tcode{/avatars} \\
Random logout on navigation & Parallel refresh requests rotated tokens twice & Single in-flight refresh in \tcode{authClient.ts} \\
Next.js dev webpack errors with avatars & \tcode{next/image} correlated with RSC issues & \tcode{ProfileAvatarImage} uses native \tcode{<img>} \\
Gemini rate limits & External API throttling under load & TTL cache + exponential backoff retry \\
Spreadsheet chunking & Tables split mid-row lost semantics & Table-aware markdown chunking \\
\bottomrule
\end{tabularx}
\end{table}

\section{Produced Technical Deliverables}

The project deliverables include:
\begin{itemize}[leftmargin=*,itemsep=0.2em]
  \item \textbf{Source code} --- monorepo at \texttt{c:\textbackslash dev\textbackslash finalproject} with \texttt{apps/api} and \texttt{apps/web};
  \item \textbf{Database} --- 29 Prisma migrations, seed script with admin/manager/employee personas;
  \item \textbf{Documentation} --- README, architecture guide, \textbf{54 PlantUML diagrams} (sequence, use-case, and sprint class flows in Chapter~\ref{ch:architecture}, full index in Appendix~\ref{annex:diagrams});
  \item \textbf{Tests} --- 10 Vitest integration suites (74 tests);
  \item \textbf{Operations} --- Docker Compose, RAG eval CLI, reindex and feedback export scripts;
  \item \textbf{UI} --- dashboard, document library, Ask interface, admin and manager consoles.
\end{itemize}

Screenshots of the running application (dashboard, Ask with sources, profile photo, admin users) should be inserted in the oral defense slide deck; the report references interfaces analytically rather than through exhaustive screenshot dumps, per ESPRIM guidance.

\section{Chapter Conclusion}

This chapter described the concrete realization of the platform across five
sprint increments, from authentication and document ingest through hybrid RAG,
governance, and operational readiness. It documented implementation constraints
and technical resolutions, showing how architectural choices were translated into
operational code, testable workflows, and reproducible deliverables.
Chapter~\ref{ch:validation} presents the validation results against the KPIs
defined in Table~\ref{tab:kpis}.
